\subsection{A sphere corner with one core orbit or at most four states}

\begin{proposition}[One core orbit for the five-element sphere corner]
\label{C:P5-one-core-bisets}
Let $P=\{1\}\sqcup(\{0,1\}\times\{0,1\})$ be the
five-element sphere monoid, with minimal band $K$, and let
$S$ be a finite nonempty $(P,N)$-biset.  Suppose that $KS$
is one two-point left $K$-orbit and that $P$ has no fixed
state in $S$.  If the collage $X$ has
\[
 \dim H^2(X;\mathbb Q)=1,\qquad H_3(X;\mathbb Q)=0,
\]
its degree-two generator has square zero.
In particular $X$ cannot have type $\mathbb{CP}^2$.
There is no bound on $|S|$ or $|N|$.
\end{proposition}
\begin{proof}
Label the core $U=\{u_0,u_1\}$ so that every band actor
$k_{ij}$ sends every state to $u_i$.  This follows from
idempotence exactly as in Proposition~\ref{C:P5-three-state-bisets}.
Commutation forces $N$ to fix both core points.  Every other
state is sent by a right actor either to another noncore
state or to one of these two points.

Put $A=\mathbb Q(S\setminus U)$, with the partial right
$N$-action that sends a basis vector to zero when its state
enters $U$.  The core of the left action category
$P\ltimes S$ is contractible.  For each noncore state there
are two cones on that core, one for each band column, with
their apex equal to that state.  Collapsing the core therefore
gives a graph with two parallel edges for each noncore state,
all sharing the core vertex.  This graph model is natural
for the right $N$-action: an actor carries a column-labeled
edge to the corresponding edge of the image state, or
collapses it to the core vertex.

For clarity, there is an explicit map from the quotient
of the actual nerve by its core to this graph.  On a
column-labeled cone send a point to the same cone-height
coordinate on its corresponding edge; equivalently this
is the barycentric weight of its apex.  The core is
contractible, so this map is a weak equivalence.  It is
$N$-equivariant: an actor preserves the height and column
label when its apex stays outside the core, and otherwise
collapses the whole cone into the core.  The graph chain complex is
$A[1]$ plus the contractible pair that sends a chosen first
edge to its noncore vertex minus the core vertex.  The
chosen edge and vertex pairs transform by the same partial
action, so this is a splitting of right $\mathbb QN$ chain
complexes.

Consequently the edge double bar has, compatibly with its
projection and core-state section to $BN$,
\[
 H_n(E_S;\mathbb Q)=H_n(BN;\mathbb Q)
       \oplus\operatorname{Tor}_{n-1}^{\mathbb QN}(A,\mathbb Q)
 \qquad(n\geq1).
\]
Its section maps constantly to $BP$.  Thus the inclusion
$BN\to X$ is null-homotopic, and cancellation of the $BN$
summands in Mayer--Vietoris gives
\[
 0\longrightarrow H_3(X;\mathbb Q)
 \longrightarrow T_1\longrightarrow\mathbb Q
 \longrightarrow H_2(X;\mathbb Q)
 \longrightarrow T_0\longrightarrow0,
 \qquad T_j=\operatorname{Tor}^{\mathbb QN}_j(A,\mathbb Q).
\]

If $H^2(X;\mathbb Q)\to H^2(BP;\mathbb Q)$ is zero,
the degree-two classes restrict to zero on both vertices.
They are relative classes, so their products vanish on the
suspension quotient by the vertex spaces.  In the remaining
case the dual map $H_2(BP;\mathbb Q)\to H_2(X;\mathbb Q)$
is an isomorphism.  The displayed sequence and the assumed
vanishing of $H_3(X;\mathbb Q)$ imply $T_0=T_1=0$.

There is an exact sequence of right $\mathbb QN$-modules
\[
 0\longrightarrow\mathbb Q(u_1-u_0)
 \longrightarrow\widetilde{\mathbb QS}
 \longrightarrow A\longrightarrow0.
\]
Its first term is the trivial module.  Rational duality
of a projective resolution of the right module $A$ gives
\[
 \operatorname{Ext}^1_{(\mathbb QN)^{\mathrm{op}}}(A,\mathbb Q)
   \cong D\operatorname{Tor}^{\mathbb QN}_1(A,\mathbb Q)=0.
\]
The sequence therefore splits.  A retraction onto its first
term defines an $N$-invariant function $q:S\to\mathbb Q$
with $q(u_i)=i$, by applying it to $s-u_0$.
Compose $q$ with the fixed function
$\rho:\mathbb Q\to\{0,1\}$ that is one at $1$ and zero
at every other rational number.  Sending $s$ to
$u_{\rho(q(s))}$ is an actual $N$-equivariant retraction
onto $U$.  It is also $P$-equivariant, since every band
actor is constant onto its prescribed $u_i$ on both sets.

The core collage is equivalent to $BP\simeq S^2$: its edge
bar is the product of a contractible core action category
with $BN$, and projection to $BN$ is an equivalence.
Thus $X$ retracts onto $S^2$.  The pullback of its sphere
class spans the one-dimensional $H^2(X;\mathbb Q)$ and
has square zero, as required.
\end{proof}

\begin{proposition}[Four states opposite the sphere corner]
\label{C:P5-four-state-bisets}
Let $S$ be a nonempty $(P,N)$-biset with $|S|\leq4$.
If $H^2(BN;\mathbb Q)=0$, its collage cannot have type
$\mathbb{CP}^2$.  When the left $P$-action has no fixed
state, the same conclusion holds without that cohomological
hypothesis on $N$.
\end{proposition}
\begin{proof}
The cases of at most three states were proved in
Proposition~\ref{C:P5-three-state-bisets}.  A $P$-fixed
state gives a section $BP\to E_S$; with the stated
vanishing at the other vertex, every degree-two collage
class is relative and its square is zero.

If there is no fixed state, every left $K$-orbit in $KS$
has two points.  One orbit is covered by
Proposition~\ref{C:P5-one-core-bisets}.  The only other
possibility at four states is $S=KS$ with two complete
orbits.  Write $S=O\times\{u_0,u_1\}$, $|O|=2$.
The right action preserves the row coordinate and acts
only on $O$.

The projection $P\ltimes S\to O$ has an $N$-equivariant
section choosing row zero.  The arrows $k_{i0}$ give an
$N$-equivariant natural transformation from that section
followed by projection to the identity.  The map of the
edge double bar to $BP$ is therefore coherently
$N$-equivariantly null-homotopic, while its map to $BN$
is identified with the action-bar projection for $O$.
The collage consequently has type
\[
 BP\ \vee\ B\mathcal C(N,O).
\]
Every positive rational product on the second summand is
zero by the two-state finite-line lemma, and the first
summand is a sphere.  The wedge has zero degree-two
products as well, completing the proof.
\end{proof}

\begin{corollary}\label{C:sixteen-P5-opposite-seven}
The sixteen-arrow profile $(5,7,4)$ needs no separate
enumeration.  If its seven-element vertex has zero
rational degree-two cohomology, apply
Proposition~\ref{C:P5-four-state-bisets}; otherwise its
opposite is in the complete $(7,5,4)$ catalog.
Any already computed $(5,7,4)$ records are independent
controls and need not be a premise of the exclusion.
\end{corollary}
\begin{proof}
Use the unique positive-degree-two five-element monoid
$P\cong P^{\mathrm{op}}$, and take the opposite category
in the second case.
\end{proof}
